% =====================================================================
%  1bat-ud1-16.tex  —  HORA 16: Prova escrita de la Unitat 1
%  Estructura: enunciat de la prova + \newpage + solucionari per al docent.
% =====================================================================
\horatitol{Sessió 16 — Prova escrita de la Unitat 1}%
{Prova individual: quatre blocs alineats amb els criteris C1–C4. El solucionari per al professorat és a la pàgina següent, separat.}

% --- Capçalera de la prova (dades de l'alumne) -----------------------
\begin{tcolorbox}[colback=white, colframe=udnavy, boxrule=0.8pt, arc=2pt,
  left=8pt, right=8pt, top=6pt, bottom=6pt]
\textbf{Prova · Unitat 1 — Nombres reals, inequacions i intervals}\\[4pt]
Nom i cognoms: \rule{6.5cm}{0.4pt} \quad Grup: \rule{1.6cm}{0.4pt} \quad Data: \rule{2.2cm}{0.4pt}\\[6pt]
\small Temps: $45$ minuts. Es valora el \textbf{resultat}, però també la
\textbf{justificació dels passos}, l'ús correcte de la \textbf{notació d'intervals},
la \textbf{comprovació} de les solucions i la \textbf{valoració} de l'ordre de
magnitud. Calculadora científica permesa. Mostra tots els càlculs.
\end{tcolorbox}

\vspace{0.4em}

% =====================================================================
\subsection*{Bloc 1 (C1) — Inequacions \hfill \normalfont\itshape ($\sim 30\%$)}
\addcontentsline{toc}{subsection}{Bloc 1 (C1) — Inequacions}

\begin{enumerate}[label=\textbf{\arabic*.}, leftmargin=2em, itemsep=8pt]
  \item Resol les inequacions següents, escriu la solució en forma d'\textbf{interval}
        i \textbf{comprova} la primera amb un valor:
    \begin{enumerate}[label=\alph*), leftmargin=1.6em, topsep=2pt]
      \item $4-2x\ge 10$ \hfill \item $3(x-2)\le 5x+4$
    \end{enumerate}
\end{enumerate}

% =====================================================================
\subsection*{Bloc 2 (C2) — Intervals a partir d'un cas real \hfill \normalfont\itshape ($\sim 25\%$)}
\addcontentsline{toc}{subsection}{Bloc 2 (C2) — Intervals}

\begin{enumerate}[label=\textbf{\arabic*.}, leftmargin=2em, itemsep=8pt, start=2]
  \item Una ajuda d'emancipació és per a joves «de $18$ anys fins a menys de $26$».
        Escriu la condició com a \textbf{inequació doble} i com a \textbf{interval},
        indica com és cada extrem (obert/tancat) i representa-la a la \textbf{recta real}.
  \item El bo social elèctric ordinari exigeix una renda que \textbf{no superi}
        $\num{12600}\eur$. Escriu la condició com a inequació i com a interval. Una
        família amb renda \emph{exactament} $\num{12600}\eur$ hi té dret? Per què?
\end{enumerate}

% =====================================================================
\subsection*{Bloc 3 (C3) — Sistema a partir d'una convocatòria \hfill \normalfont\itshape ($\sim 20\%$)}
\addcontentsline{toc}{subsection}{Bloc 3 (C3) — Sistema}

\begin{enumerate}[label=\textbf{\arabic*.}, leftmargin=2em, itemsep=8pt, start=4]
  \item Un ajut al lloguer jove demana complir alhora: \textbf{edat} dins de
        $[18,35]$ \emph{i} \textbf{renda} anual $\le\num{24000}\eur$. Decideix,
        justificant-ho, si hi tenen dret:
    \begin{enumerate}[label=\alph*), leftmargin=1.6em, topsep=2pt]
      \item una persona de $28$ anys amb renda $\num{22000}\eur$
      \item una persona de $37$ anys amb renda $\num{15000}\eur$
    \end{enumerate}
\end{enumerate}

% =====================================================================
\subsection*{Bloc 4 (C4) — Notació científica \hfill \normalfont\itshape ($\sim 25\%$)}
\addcontentsline{toc}{subsection}{Bloc 4 (C4) — Notació científica}

\begin{enumerate}[label=\textbf{\arabic*.}, leftmargin=2em, itemsep=8pt, start=5]
  \item (a) Escriu en notació científica: $\num{3326000000}$ i $\num{1950000}$.\\
        (b) Un pressupost de $3,326\times 10^{9}\eur$ es reparteix entre
        $1,95\times 10^{6}$ persones. Calcula els euros per persona i \textbf{valora}
        si l'ordre de magnitud té sentit.
\end{enumerate}

% =====================================================================
\newpage
% =====================================================================
\fullsolucions{Prova UD1 — per al professorat}
\begin{solucions}

\noindent\textbf{Bloc 1 (C1)}
\begin{sollist}
  \item (a) $4-2x\ge 10 \Rightarrow -2x\ge 6$; dividim per $-2$ i \textbf{girem}:
        $x\le -3$, interval $(-\infty,-3]$. Comprovació amb $x=-4$:
        $4-2(-4)=12\ge 10$, cert. \quad
        (b) $3x-6\le 5x+4 \Rightarrow -2x\le 10$; girem: $x\ge -5$, interval
        $[-5,+\infty)$.
\end{sollist}

\noindent\textbf{Bloc 2 (C2)}
\begin{sollist}\setcounter{enumi}{1}
  \item «de $18$ fins a menys de $26$»: $18\le x<26$, interval $[18,26)$; el $18$
        tancat (hi entra), el $26$ obert (no hi entra). A la recta: punt ple a $18$,
        punt buit a $26$, segment ressaltat entremig.
  \item $\text{renda}\le\num{12600}$, interval $(-\infty,\,\num{12600}]$. Una família
        amb renda exactament $\num{12600}\eur$ \textbf{sí} que hi té dret, perquè la
        condició és $\le$ («no superar»): el límit hi entra (claudàtor).
\end{sollist}

\noindent\textbf{Bloc 3 (C3)}
\begin{sollist}\setcounter{enumi}{3}
  \item (a) Edat $28\in[18,35]$, compleix; renda $\num{22000}\le\num{24000}$,
        compleix. Compleix \textbf{totes dues} $\rightarrow$ \textbf{té dret}. \quad
        (b) Edat $37\notin[18,35]$ ($37>35$): falla l'edat. Encara que la renda
        ($\num{15000}\le\num{24000}$) compleixi, com que cal complir-les totes,
        \textbf{no té dret}.
\end{sollist}

\noindent\textbf{Bloc 4 (C4)}
\begin{sollist}\setcounter{enumi}{4}
  \item (a) $\num{3326000000}=3,326\times 10^{9}$ i
        $\num{1950000}=1,95\times 10^{6}$. \quad
        (b) $\dfrac{3,326\times 10^{9}}{1,95\times 10^{6}}=\dfrac{3,326}{1,95}\times 10^{3}
        \approx 1,71\times 10^{3}\eur\approx\num{1710}\eur$ per persona. L'ordre de
        magnitud (milers d'euros) \textbf{té sentit} per a un ajut anual per persona.
\end{sollist}

\end{solucions}

\noindent\small\textit{Orientació de qualificació (rúbrica): Bloc 1 $\sim 30\%$,
Bloc 2 $\sim 25\%$, Bloc 3 $\sim 20\%$, Bloc 4 $\sim 25\%$. A cada bloc es reparteix
entre el procediment correcte, la notació (intervals, extrems, notació científica) i
la justificació/comprovació. Llindar de superació: $4,0/10$; en cas contrari,
recuperació trimestral.}
